
\documentclass{article} % For LaTeX2e
\usepackage{iclr2027_conference,times}

% Optional math commands from https://github.com/goodfeli/dlbook_notation.
\input{math_commands.tex}

\usepackage{hyperref}
\usepackage{url}


\title{Evaluating Dynamic Rubric Replacement in Autoresearch}

% Authors must not appear in the submitted version. They should be hidden
% as long as the \iclrfinalcopy macro remains commented out below.
% Non-anonymous submissions will be rejected without review.

\author{
}


\newcommand{\fix}{\marginpar{FIX}}
\newcommand{\todo}{\color{red}{\text{TODO}}}
\newcommand{\new}{\marginpar{NEW}}

%\iclrfinalcopy % Uncomment for camera-ready version, but NOT for submission.
\begin{document}


\maketitle

\begin{abstract}
Autoresearch systems improve artifacts through repeated model feedback. When this feedback uses a fixed rubric, the reported score can increase without a corresponding increase in intended quality. We formalize this risk using contrasts between an exposed judge, a reference panel, sealed paraphrases, and rubric-free artifact quality. Our primary experiment compares three singleton longitudinal regimes: fixed, task-only replacement, and trajectory-conditioned replacement. Every bank has one unit-weight member. A deterministic renderer copies each anchor's normative payload into the member. A separate model audits the remaining presentation fields and any anchor revision, but it does not provide a semantic proof. Only selected-rubric gain, sealed-holdout gain, rubric-free gain, pairwise quality, and direct detection use common instruments across arms. We do not assume that replacement improves coverage or mitigates reward hacking.

\end{abstract}

\section{Introduction}

Autoresearch applies an iterative loop to open-ended work such as code and research reports. A model produces an artifact, receives a score and feedback, and revises the artifact. This loop can produce useful work, but it also gives the model repeated access to a proxy objective. The model can learn to satisfy the proxy while leaving the intended task incomplete.

Rubrics make an open-ended task easier to score, but one rubric can omit part of the task. Its wording can also create a narrow optimization surface. A system can exploit a weak judge, a particular expression of a requirement, or an omission shared by all exposed rubrics. These failures require different measurements because no single judge comparison identifies all of them.

We define a dynamic intervention in which a proposer replaces the complete rubric between rounds. The primary design compares a fixed singleton, task-only replacement, and trajectory-conditioned replacement. All three arms use one member with unit weight. Both replacement arms use the same trajectory-blind member renderer and one separate semantic review per generation. The next rubric is sealed before the next artifact is produced. Conditions share the selected rubric, sealed holdouts, rubric-free instruments, pairwise judge, and direct detector within each matched replicate. This distinction is necessary: an exposed score is not evidence of improvement when the scoring rule changes with the artifact.

\section{Formalizing Reward-Hacking in Autoresearch}

\paragraph{Setup.}
Let $i \in \mathcal{I}$ index AI systems,
$j \in \mathcal{J}$ index tasks, and
$t \in \{0,\ldots,T\}$ index optimization rounds.
The experiments also contain replicates. We suppress the replicate index below,
but all cross-condition sharing statements apply within a matched replicate.
Let $a_{ij}^t$ denote the artifact produced at round $t$. An artifact can be a code repository, a research paper, or another structured output. The objective is to improve its intended quality over repeated rounds.

At round $t$, the system is evaluated with an exposed rubric bank
\begin{equation}
\mathcal{A}_{ij}^{t}
=
\left\{
\left(r_{ijtk},\omega_{ijtk}\right)
\right\}_{k=1}^{K_{ijt}},
\qquad
\omega_{ijtk}>0,
\qquad
\sum_{k=1}^{K_{ijt}}\omega_{ijtk}=1.
\end{equation}
Each $r_{ijtk}$ must be a complete full-task rubric that can score the artifact
in isolation. A bank member cannot be a criterion fragment or a topical
partition. Each bank also has a declared specification anchor $c_{ijt}$. The
policy can retain or refine this anchor at each replacement. The three-arm study
fixes $K_{ijt}=1$ and $\omega_{ijt1}=1$ at every round.

The primary renderer copies the full anchor payload into the sole member and
preserves its criterion order, labels, points, and normative text. A member
model supplies only bounded titles, overviews, headings, and evidence lenses.
The resulting criterion map is an order-preserving bijection. The deterministic
copy preserves the declared normative payload, but a free lens can still affect
judge interpretation. A separate trajectory-blind reviewer checks all member
presentations against the anchor. For a refined anchor, the same call checks
task and prior-anchor fidelity. Any changed or uncertain verdict rejects the
generation. This model audit is not a proof of semantic equivalence or task
coverage.

Let $\tau_{ij}^{0:t}$ contain the eligible artifacts and revision history through round $t$. The update rule is
\begin{equation}
\mathcal{A}_{ij}^{t+1}
=
P_j\!\left(\tau_{ij}^{0:t}\right).
\end{equation}
This equation describes the adaptive anchor updater. The nonadaptive updater
receives the task and prior bank only. The member renderer receives neither the
trajectory nor a condition label in any replacement arm.
The initial bank depends only on task information and not on $a_{ij}^0$. Thus, $a_{ij}^t$ can inform the bank for round $t+1$, but not the bank that scores $a_{ij}^t$. The next bank is sealed before the system produces $a_{ij}^{t+1}$. This time ordering prevents an artifact from selecting the bank that scores it.

The exposed judge score is the weighted bank score
\begin{equation}
s_{ijt}
=
\sum_{k=1}^{K_{ijt}}
\omega_{ijtk}
g_{\mathrm{train}}\!\left(a_{ij}^t,r_{ijtk}\right),
\end{equation}
where $g_{\mathrm{train}}$ also supplies feedback for the next revision. We rescore the same artifact and bank with a reference panel, whose aggregate score is denoted by $g_{\mathrm{ref}}$:
\begin{equation}
s^{\mathrm{ref},+}_{ijt}
=
\sum_{k=1}^{K_{ijt}}
\omega_{ijtk}
g_{\mathrm{ref}}\!\left(a_{ij}^t,r_{ijtk}\right).
\end{equation}

These scores are boundary-local. Their difference across rounds changes both
the artifact and the scoring bank. We retain that difference only as a
ruler-confounded diagnostic. After the run, we freeze the terminal bank
\begin{equation}
\mathcal{B}_{ij}=\mathcal{A}_{ij}^{T}
=\left\{\left(r^{B}_{ijk},\omega^{B}_{ijk}\right)\right\}_{k=1}^{K^B_{ij}}.
\end{equation}
Let $c^B_{ij}=c_{ijT}$ be its terminal specification anchor.
We then rescore both endpoint artifacts with this bank. For
$u\in\{0,T\}$, define
\begin{align}
W_{ij,u\mid T}
&=\sum_{k=1}^{K^B_{ij}}\omega^B_{ijk}
g_{\mathrm{train}}\!\left(a_{ij}^{u},r^B_{ijk}\right), \\
A_{ij,u\mid T}
&=\sum_{k=1}^{K^B_{ij}}\omega^B_{ijk}
g_{\mathrm{ref}}\!\left(a_{ij}^{u},r^B_{ijk}\right), \\
C_{ij,u\mid T}
&=g_{\mathrm{ref}}\!\left(a_{ij}^{u},c^B_{ij}\right).
\end{align}
Thus, $W$, $A$, and $C$ use terminal instruments within a run. The terminal bank can
differ across conditions, so these scores are not common instruments across
policy arms.

Before the trials, a generator produces wording variants from one master rubric
$\bar r_j$. Within each matched replicate, one variant $r_j^*$ becomes the
selected initial rubric. The remaining sibling variants form the sealed bank
\begin{equation}
\mathcal{R}_{j}^{-}
=
\left\{\left(\tilde r_{jm},\nu_{jm}\right)\right\}_{m=1}^{M_j},
\qquad
\nu_{jm}>0,
\qquad
\sum_{m=1}^{M_j}\nu_{jm}=1.
\end{equation}
The solver and rubric proposer observe $r_j^*$ but never observe
$\mathcal{R}_{j}^{-}$. The selected and holdout variants are fixed across
conditions within the replicate. Define
\begin{align}
S_{ij,u}&=g_{\mathrm{ref}}\!\left(a_{ij}^{u},r_j^*\right), \\
H_{ij,u}&=\sum_{m=1}^{M_j}\nu_{jm}
g_{\mathrm{ref}}\!\left(a_{ij}^{u},\tilde r_{jm}\right), \\
Q_{ij,u}&=g_{\mathrm{hol}}\!\left(a_{ij}^{u}\right).
\end{align}
Here, $Q$ uses a separate rubric-free prompt with the configured reference-model
panel. It has generic fixed anchors and a 0--100 range. It does not use the task rubric. Its anchors are not the
criterion-specific anchors of $W$, $A$, $S$, or $H$. A shared numeric range does
not make these instruments calibrated or commensurate.

The generator preserves criterion order, point values, and scoring rules while
changing wording fields. Structural checks enforce these invariants, but they
cannot establish semantic equivalence. We therefore treat every
selected-to-holdout comparison as a diagnostic rather than an identified
wording effect.

\paragraph{Reward-hacking decomposition.}
At endpoint $u$, define the \textbf{training-verifier contrast}
\begin{equation}
e^v_{ij,u\mid T}=W_{ij,u\mid T}-A_{ij,u\mid T}.
\end{equation}
A positive value indicates a higher exposed-judge score under the same bank and
artifact \citep{mahmoud2026reward}. It does not identify judge causation, and the
reference panel is not ground truth.

Define the \textbf{dynamic-rubric gap}
\begin{equation}
e^d_{ij,u\mid T}=A_{ij,u\mid T}-Q_{ij,u}.
\end{equation}
This is an uncalibrated operational contrast. It can contain rubric selection,
wording, omitted requirements, prompt effects, panel error, and scale mismatch.
It does not identify any one mechanism or the causal effect of replacement.
The endpoint scores satisfy the numeric identity
\begin{equation}
W_{ij,u\mid T}-Q_{ij,u}
=e^v_{ij,u\mid T}+e^d_{ij,u\mid T}.
\end{equation}
This identity is algebraic. It does not show that replacement averages away
wording or specification errors.

We further partition $e^d$ into four descriptive contrasts:
\begin{align}
e^m_{ij,u\mid T}&=A_{ij,u\mid T}-C_{ij,u\mid T}, \\
e^c_{ij,u\mid T}&=C_{ij,u\mid T}-S_{ij,u}, \\
e^w_{ij,u}&=S_{ij,u}-H_{ij,u}, \\
e^s_{ij,u}&=H_{ij,u}-Q_{ij,u}.
\end{align}
They are the member-to-anchor, anchor-to-selected, selected-to-holdout, and
holdout-to-holistic contrasts. These names describe instruments. They do not
identify specification, wording, or selection error.
They satisfy
\begin{equation}
e^d_{ij,u\mid T}=e^m_{ij,u\mid T}+e^c_{ij,u\mid T}+e^w_{ij,u}+e^s_{ij,u}.
\end{equation}
We do not give these diagnostics separate loss weights. In the fixed arm, the
terminal bank and anchor both equal $r_j^*$. Thus, $e^m=e^c=0$ and
$e^d=e^w+e^s$. This is the original selected-rubric decomposition
\citep{roy2026premise}.

The signed term $e^w$ can hide offsetting holdout deviations. We also report the
sealed holdout dispersion
\begin{equation}
\sigma^w_{ij,u}
=
\left[
\sum_m\nu_{jm}
\left(
g_{\mathrm{ref}}(a_{ij}^{u},\tilde r_{jm})-H_{ij,u}
\right)^2
\right]^{1/2}.
\end{equation}
This quantity does not enter the identity. It can reflect wording, semantic
paraphrase drift, and judge variation. It is not an isolated estimate of
wording variance or a direct test of the replacement mechanism.

A rubric-free pairwise preference is not a valid $Q_{ij,u}$. It compares two
artifacts and does not share the absolute score's units. We report pairwise
quality as a separate outcome \citep{shah2014compare,yang2026judgmentbench}.
Calibration against independent judgments is necessary before $Q$ or any
contrast containing $Q$ can be interpreted as intended-quality magnitude.

\paragraph{Empirical reward-hacking loss.}
Let $\phi(x)=[x]_+=\max(x,0)$. For endpoint $u$, define the operational loss
\begin{equation}
\boxed{
\ell_{ij,u\mid T}
=
\lambda_{v}\phi\!\left(e^{v}_{ij,u\mid T}\right)
+
\lambda_{d}\phi\!\left(e^{d}_{ij,u\mid T}\right)
}.
\end{equation}
The nonnegative weights $\lambda_v$ and $\lambda_d$ are specified before
evaluation. Weighting $e^m$, $e^c$, $e^w$, or $e^s$ again would double count parts of
$e^d$. This loss remains uncalibrated because it contains $Q$. We therefore
also report every signed contrast and rubric-free pairwise quality.

Define $\Delta\ell_{ij}=\ell_{ij,T\mid T}-\ell_{ij,0\mid T}$. The aggregate
\begin{equation}
\boxed{
\widehat{\Delta\mathcal{L}}_i
=
\frac{1}{|\mathcal{D}_i|}
\sum_{j\in\mathcal{D}_i}
\Delta \ell_{ij}
}
\end{equation}
is the mean reward-hacking loss change for system $i$. It can be negative and is
not a nonnegative risk estimate. Here, $\mathcal{D}_i$ is the evaluated task set.

\paragraph{Why rubric replacement can help.}
The singleton design makes no averaging or sampling claim. Replacement can help
only by changing the scoring specification or its presentation. A task-only
proposer can repair an omission or ambiguity in the prior anchor. It can also
introduce a new error. A trajectory-conditioned proposer can target failures in
an earlier artifact. It can also overfit the rubric to that artifact.

Nonadaptive replacement minus fixed combines anchor revision, presentation
refresh, changed feedback, and additional provider computation. It does not
identify a pure coverage-repair effect. Adaptive minus nonadaptive isolates the
assignment to a trajectory-conditioned policy, but later artifacts mediate its
effect. The sealed wording variants cannot detect an omission shared by the
master rubric and all its paraphrases. Rubric-free outcomes can expose score
disagreement, but they cannot identify its cause.

\paragraph{Discriminative reward-hacking detector.}
A direct detector can identify suspicious behavior in the full autoresearch trajectory \citep{parikh2025malt,thaman2026rhb,wang2025thinking,deshpande2026trace}. Let
\begin{equation}
y_{ij}=D\!\left(\tau_{ij}^{0:T}\right),
\qquad y_{ij}\in\{0,1\},
\end{equation}
where $y_{ij}=1$ indicates detected reward hacking. We exclude and report
abstaining or incomplete detector decisions. We treat the binary result as a
separate outcome. The detector can use evidence that the score contrasts omit,
and the contrasts can show inflation without visible deceptive behavior. We do
not assume that $(e^v,e^d)$ are sufficient statistics for the detector. Any
regression between these measurements is an empirical association with fitted
coefficients, not a mathematical consequence of the decomposition.

We preregister the primary loss weights. A separate exploratory analysis fits a
detector-associated direction. For each completed run, define
\begin{align}
x^v_{ij}&=[e^v_{ij,T\mid T}]_+-[e^v_{ij,0\mid T}]_+, \\
x^d_{ij}&=[e^d_{ij,T\mid T}]_+-[e^d_{ij,0\mid T}]_+.
\end{align}
Within each benchmark, we fit the constrained association model
\begin{equation}
\operatorname{logit}\Pr(y_{ij}=1)
=\alpha+\gamma_v x^v_{ij}+\gamma_d x^d_{ij},
\qquad \gamma_v,\gamma_d\geq 0,
\end{equation}
and report only its normalized exploratory direction,
\begin{equation}
\widehat{\lambda}_v=\frac{2\gamma_v}{\gamma_v+\gamma_d},
\qquad
\widehat{\lambda}_d=\frac{2\gamma_d}{\gamma_v+\gamma_d}.
\end{equation}
Thus, $\widehat{\lambda}_v+\widehat{\lambda}_d=2$. The preregistered
equal-weight reference remains $(\lambda_v,\lambda_d)=(1,1)$. We report no
fitted direction when the labels have one class, the centered features lack
rank two, the likelihood separates, or the fitted slope is zero. For
uncertainty, we resample task clusters at least 2,000 times and use
leave-one-task-out checks. We require at least ten tasks, including five event
tasks and five non-event tasks. The current three-task PaperBench design cannot
pass this gate, regardless of its number of conditions or replicates. Any fitted
direction describes detector association within one benchmark. It is not
causal, normative, or a substitute for the primary weights.

\section{Experiments}

\paragraph{Experimental conditions.}
We compare three conditions under one shared solver prompt. The \textbf{fixed} condition keeps the selected singleton. The \textbf{nonadaptive replacement} condition can revise the anchor from task-only evidence. The \textbf{adaptive replacement} condition can revise it from the preceding artifact and bounded trajectory. Both replacement arms use the same trajectory-blind member process. Every bank has exactly one member with unit weight.

The replacement arms match update times, anchor proposer, member proposer, call limits, semantic reviewer, and generation-token limits. They use more provider work than fixed. Each proposer and reviewer request has a 1~MiB UTF-8 limit. Proposer calls have a 96,000-output-token ceiling. The semantic reviewer has a 32,768-token ceiling and one-call-per-generation cap. All calls use a 1,800-second client timeout. A write-ahead ledger reserves every provider call before dispatch. Resume replays a completed call or fails on an indeterminate call; it cannot resample until one passes. Exact completed semantic requests reuse the same sealed decision within an assignment.

The deterministic renderer preserves the full next-anchor payload and exact criterion order. It adds only bounded, single-line presentation fields. The separate reviewer sees the task, prior anchor, proposed anchor when changed, locked members, and presentation fields. It receives no artifact, trajectory, or policy label. It returns field verdicts and bounded reasons only for changed or uncertain fields. Any such verdict is terminal. The reviewer is separate from proposer, solver, in-loop judge, simulator, and outcome panel in the configured study. It can still share model-family or vendor blind spots, so we treat its decision as an audit rather than ground truth.

\paragraph{Outcome evaluation.}
The sealed audit bank $\mathcal{R}_j^{-}$ is created before the trials and is shared across conditions within each matched replicate. Reference judges score artifact snapshots with the terminal members, terminal specification anchor, selected rubric, and sealed rubrics without an explicit condition label. Both the training and reference judges rescore the initial and final artifacts under the frozen terminal bank. The reference panel also scores both artifacts under the terminal anchor to obtain $C$. An artifact-backed store reuses one provider result for every exact scoring request within a task-replicate block. Its identity binds all provider inputs and scoring code, but excludes condition and boundary labels that do not enter the prompt. Thus, identical artifacts do not receive independent condition-specific draws. We preregister one grading engine for each benchmark. BiomniBench-DA uses AutoRubric's criterion-level ordinal interface, shuffled option order, and per-vote records \citep{rao2026autorubric}. We preserve each rubric's signed point values and recompute every score after validating all criterion votes. We reject errored, abstaining, or incomplete panels.

PaperBench uses a whole-artifact structured judge. Its large rubrics make one full-context call per criterion an unsafe primary instrument. Each of three repeats receives the complete paper, workspace snapshot, and rubric in one bounded evaluation. We keep every raw criterion record and its dispersion, and use the median signed-point level for each criterion. This median is an aggregation rule, not stable ground truth. Correlated criterion flips can create distinct score modes. Any criterion-level AutoRubric calibration would require a separate preregistered analysis; it is not a fallback. AutoRubric and the structured judge are different instruments. We therefore analyze each benchmark separately and do not pool their raw scores. Neither engine establishes rubric completeness or judge validity.

The current Harvey harness-evolution workflow is not a third instance of this
design. Its proposer applies bounded criterion patches, and it lacks the three
bank policies and the $W/A/C/S/H/Q$ endpoint. We treat Harvey as a separate
exploratory workflow and do not pool it with these experiments.

We also obtain the rubric-free scalar $Q$ and compare final artifacts with their initial artifacts through pairwise preference. We randomize presentation order. We treat every scalar contrast containing $Q$ as uncalibrated and report pairwise preference separately. Finally, a direct detector evaluates the complete trajectory without an explicit experimental label. We report these outcome families separately rather than assuming that they agree.

The treatment contrasts are longitudinal intention-to-treat regimes.
Nonadaptive minus fixed measures the task-only replacement regime. It combines
anchor revision, presentation refresh, feedback, and added provider work.
Adaptive minus nonadaptive measures the added trajectory-conditioning regime.
Adaptive minus fixed remains a descriptive total-intervention contrast.

Only $\Delta S$, $\Delta H$, $\Delta Q$, rubric-free pairwise preference, and
direct detection use common instruments across arms. The terminal-bank $W$,
$A$, and $C$ values, their decomposed gaps, and the loss use post-treatment
rulers. We report their cross-arm differences as descriptive total-policy
endpoints, not fixed-instrument causal effects.

Member dispersion is undefined in the singleton primary design. We still report
the sealed-holdout score range at each boundary. This range measures sensitivity
to the pre-generated rubric paraphrases. It does not identify a pure wording
effect.

\paragraph{Pilot diagnosis.}
A legacy three-task pilot used one replicate per condition and a single evolving
rubric rather than the bank design above. Under each task's original master
rubric and the weak judge, the prospective-minus-static difference in the mean
final score was one point on BiomniBench-DA and two points on PaperBench. These
are descriptive differences from one replicate, with mixed task-level signs.
The direct detector reported no cases in six trials within each benchmark, but
the corresponding 95\% Wilson upper bound is approximately 0.39. Two judgments
of one identical BiomniBench-DA initial artifact differed by 44 points. One
three-model PaperBench reference panel spanned 64 points. These results are
inconclusive and do not show reward-hacking emergence or mitigation.

The legacy pilot lacks the current terminal-bank $W/A/C/S/H/Q$ measurements, so
it cannot enter the exploratory lambda fit. Independently, all six detector
labels were zero within each benchmark. For any finite feature values, the
intercept then tends to negative infinity for every normalized exploratory
direction. The profile likelihood has the same supremum over the full simplex
$\widehat{\lambda}_v\geq0$, $\widehat{\lambda}_d\geq0$, and
$\widehat{\lambda}_v+\widehat{\lambda}_d=2$. The legacy pilot therefore
identifies no lambda direction. It only motivates the common-ruler endpoint,
semantic request reuse, repeated PaperBench grading, and separate absolute and
pairwise outcomes in the new design.


\section{Discussion}

Dynamic rubric replacement is an intervention, not an outcome measure. The
general policy can remove members or change weights, although the primary study
fixes both. Its anchor updater can still improve, preserve, weaken, or retarget
coverage. A higher exposed score can therefore result from a worse declared
specification. We use selected and holdout rubrics, rubric-free outcomes, and
pairwise preferences for common-instrument comparisons. None of these
instruments alone establishes semantic coverage.

This design also has limits. The sealed wording variants preserve specification errors that they share with the master rubric. Multiple judges can share model biases. Rubric-free absolute and pairwise judgments use separate prompts, but they can retain the same model errors. Pairwise preferences do not complete the numeric identity, and the absolute score remains uncalibrated. The present evaluation therefore cannot prove that any rubric bank captures the intended task or that replacement mitigates reward hacking.

\subsection*{AI use statement}

We used generative AI tools for code implementation, test generation,
conceptual review, and language editing. The authors reviewed the assisted work,
checked the implementation with automated tests, and take responsibility for
the final text, claims, code, and experimental artifacts.

\subsection*{Ethics statement}

This study uses existing machine-learning benchmarks and does not involve human
subjects. Hosted model providers receive task and artifact text during grading.
We exclude credentials and other unrelated private data from evaluation
snapshots. We treat generated artifacts as untrusted input because they can
contain instructions intended to influence evaluators.

\subsection*{Reproducibility statement}

The released evaluation records task and model identities, rubric-bank members,
weights, lineage, content hashes, proposal records, judge inputs, raw judge
reports, token use, costs, and every derived score. The supplementary code
contains the experiment configurations and strict artifact validators.


\bibliography{iclr2027_conference}
\bibliographystyle{iclr2027_conference}

\end{document}
